%% ma: L12-B08-0095
%% lop: 12
%% chuong: 2
%% bai: 8
%% dang: 12.8.29
%% loai: TN4
%% muc: NB
%% dap_an: "B"
%% nguon: "CĐ2.3.pdf giáo viên gửi 27/10/2026 (M07, Lê Minh Khôi – Nguyễn Anh Khôi, Chương II Chủ đề 3 Biểu thức tọa độ của các phép toán vectơ), Câu 3"
%% kien_thuc: "Bài 8 (biểu thức tọa độ của phép cộng, trừ vectơ, tích vectơ với một số)"
%% trang_thai: chinh-thuc
\begin{ex}
Trong không gian $Oxyz$, cho hai vectơ $\vec u=(1;4;-5)$ và $\vec v=(-2;-1;3)$. Tọa độ của vectơ $\vec k=\vec u-2\cdot\vec v$ là
\choice{$\vec k=(-5;-6;11)$}{\True $\vec k=(5;6;-11)$}{$\vec k=(5;6;1)$}{$\vec k=(1;5;-8)$}
\loigiai{$2\vec v=(-4;-2;6)$, $\vec k=(1+4;4+2;-5-6)=(5;6;-11)$.}
\end{ex}
